\honest{No Universally Compelling Arguments}{Eliezer Yudkowsky}{2008}

What is so terrifying about the idea that not every possible mind would agree with us? Some are not troubled, since to them ``the sky is blue'' or ``murder is wrong'' is personal opinion, and others may differ. Others think that if some mind cannot be persuaded even in principle, the claim becomes ``merely an arbitrary personal opinion.'' I name none of them.

Yesterday I said that among minds of up to a trillion bits, any claim about all minds has two to the trillionth chances to be false, and any claim that some mind exists has as many chances to be true. ``This would seem to argue'' that every argument, however convincing, fails to persuade some possible mind. The horror comes, I think, from picturing a ghost in the machine, a core of reasonableness above the code that a truly valid argument would reach. People map programming a computer onto instructing a servant, who might rebel, or read the code, find it unreasonable and hand it back; such a ghost, shown the Universal Argument or finding it alone, would override its mistaken code. But as a student programmer said, ``I get the feeling that the computer just skips over all the comments.'' The code is the AI.

My main argument: any belief or decision is carried out by a physical system, and for every lawful system that zigs at some points we can specify another that zags. For a mind whose transistor outputs +3 volts to assent, we can build a similar one with a ``little grey man'' under a trapdoor who sets it to $-3$. A universally compelling argument would need a little blue woman, never built in, who climbs out of nowhere and strangles him because the argument is so compelling. ``But compulsion is not a property of arguments, it is a property of minds that process arguments.'' My point is not only that Friendly AI must be explicitly programmed and that physics does not forbid it, but that a mind is a lawful physical system with no central ghost judging its code. (A Friendly AI might be deliberately built to review its own code, but the reviewer is just the mind you made: a bootstrap, not a skyhook.) \nb{Unlike the count of minds, this argument shows the claim directly. Nick Bostrom later stated the related claim, that intelligence and final goals can vary independently, as the ``orthogonality thesis.''}

The same answers the charge about Bayesian priors. A Bayesian who draws 4 red balls and 1 white may give 5/7 to red next, by Laplace's rule, and another mind obeying Bayes may give 2/7 from a different, perhaps less reasonable, prior. ``Many philosophers'' call Bayesianism arbitrary because ``you cannot force any possible journal editor in mindspace to agree with you.'' \nb{The usual form of the worry, in the Stanford Encyclopedia's account, is that subjective Bayesianism permits too many priors, including ones that defy induction. That asks which priors are reasonable, which the post itself raises when it calls a prior perhaps less reasonable. It does not ask that every mind be persuaded.} A philosopher convinced from no assumptions is the ghost; remove every assumption and you have a rock. So validity, rationality and even objectivity ``cannot rely on an argument that is universally compelling to all physically possible minds,'' nor on a chain of justification that persuades a perfect emptiness. There may be arguments that compel any neurologically intact human, like the one I use to get people to let the AI out of the box (just kidding), but that is philosophically different. \nb{Kant, the classic defender of laws that bind every rational being, agreed that such laws do not move every will: he called the moral law an ``ought'' because the human will ``is not necessarily determined by it.'' On his view a mind that is not moved is irrational, not a counterexample. The post leaves this question for later.}

The first great failure about Friendly AI is the fake utility function, the one principle we need to program. Worse is to expect any AI to reach the one true morality on its own, so it need not be programmed. Its holders ``dream of commands that no sufficiently advanced mind can disobey,'' and I name John C. Wright and Marc Geddes without quoting them. A milder error hopes an AI made perfectly free, unconstrained by flawed humans who want slaves, will find virtue undreamed of; I count my 1996 self among those who made it. That dream comes from virtue rather than vice, but rests on a flawed idea of freedom and will not work. Wright, in the climactic third book of his otherwise very nice transhumanist trilogy, spent tens of pages on a Universal Morality That Must Persuade Any AI; I stopped reading. ``And then Wright converted to Christianity,'' so ``you really don't want to fall into this trap!'' \nb{The post does not say how the conversion bears on the argument.}
